\addcontentsline{toc}{section}{Phụ lục}

Phần phụ lục này cung cấp liên kết đến các tài liệu hỗ trợ cho báo cáo, bao gồm chính sách pháp lý, khảo sát nhu cầu người dùng và các tài liệu thiết kế bổ sung.

\subsection*{A. Tài liệu pháp lý liên quan}
\addcontentsline{toc}{subsection}{A. Tài liệu pháp lý liên quan}

Các quy định pháp lý được sử dụng làm cơ sở để đánh giá tính hợp lệ của mô hình cho thuê xe máy điện, giao dịch điện tử và xử lý dữ liệu cá nhân:

\url{https://drive.google.com/file/d/17tyYIoYQgS55XRo3lMkTFauNXceZZbh2/view?usp=drive_link}

\vspace{10pt}

\subsection*{B. Kết quả khảo sát nhu cầu cộng đồng}
\addcontentsline{toc}{subsection}{B. Kết quả khảo sát nhu cầu cộng đồng}

Khảo sát được thực hiện để đánh giá mức độ quan tâm đối với mô hình thuê xe máy điện theo giờ, các tính năng người dùng mong muốn và thói quen sử dụng phương tiện di chuyển.

\begin{itemize}
    \item Form khảo sát (Google Form):
    
    \url{https://docs.google.com/forms/d/e/1FAIpQLScdU2fNqEeZQIFWnPs1rlXdPqpPLl07l_Q1lZIT_emHEnrG8g/viewform?usp=dialog}

    \item Bảng tổng hợp kết quả khảo sát (Google Sheets):  
    
    \url{https://drive.google.com/file/d/1qalHs-whOtyt8GSPV6-ShBEndhdr-MAg/view?usp=sharing}
\end{itemize}

\vspace{10pt}

\subsection*{C. Tài liệu thiết kế bổ sung}
\addcontentsline{toc}{subsection}{C. Tài liệu thiết kế bổ sung}

Bao gồm các tài liệu không đưa đầy đủ vào phần nội dung chính của báo cáo.

\begin{itemize}
    \item Prototype giao diện trên Figma:  
    
    \url{https://www.figma.com/proto/PM7XgY1isY4DmazIL1I40o/P2P-Electric-Scooter-Sharing-App?node-id=0-1&t=xyMUibPTbahANB60-1}

    % \item Sơ đồ ERD phiên bản đầy đủ:  
    % \url{https://drive.google.com/your-erd-link}

    % \item Tài liệu API (OpenAPI/Swagger):  
    % \url{https://your-api-docs-link}
\end{itemize}

\vspace{10pt}

\subsection*{D. Mã nguồn hệ thống}
\addcontentsline{toc}{subsection}{D. Mã nguồn hệ thống}

Toàn bộ mã nguồn của hệ thống được quản lý tập trung trong một tổ chức GitHub, bao gồm mã nguồn phía máy chủ (backend), ứng dụng người dùng (mobile app) và giao diện quản trị (admin web).  
Mã nguồn được công bố nhằm phục vụ mục đích tham khảo, đánh giá kỹ thuật và phát triển mở rộng trong tương lai.

\url{https://github.com/DACN-P2P-Electric-Motorbike}
